% TOP_PROBLEM: {"id": 2021, "title": "Erdős Problem #271", "category_id": 1, "category": {"id": 1, "name": "number_theory", "display_name": "Number Theory", "description": "Properties of integers, prime numbers, Diophantine equations.", "slug": "number-theory", "order_index": 1, "created_at": "2026-07-31T15:26:25.670Z"}, "status": "open"}
% FIRST_POSTED: null
% ENTRY_KIND: statement_only
% Imported from: unsolvedmath_additional_500.tex; UnsolvedMath ID 2021
% !TeX program = lualatex
% Compile twice: lualatex 2021.tex
% Self-contained: no external bibliography, images, or \input files.
% Numeric section labels and visible numbers are original UnsolvedMath IDs.
\documentclass[11pt,a4paper]{article}
\usepackage[margin=24mm,headheight=14pt]{geometry}
\usepackage{fontspec}
\setmainfont{Latin Modern Roman}
\setsansfont{Latin Modern Sans}
\setmonofont{Latin Modern Mono}
\usepackage{amsmath,amssymb,mathtools}
\usepackage{unicode-math}
\setmathfont{Latin Modern Math}
\usepackage{microtype}
\usepackage{enumitem,xurl,needspace}
\usepackage[dvipsnames]{xcolor}
\usepackage{hyperref}
\hypersetup{unicode=true,colorlinks=true,linkcolor=MidnightBlue,urlcolor=MidnightBlue,
 pdftitle={UnsolvedMath Problem 2021},pdfauthor={Source-annotated compilation},bookmarksnumbered=true}
\usepackage{bookmark,tocloft,titlesec,fancyhdr}
\setlength{\cftsecnumwidth}{4.2em}
\cftsetpnumwidth{2.5em}
\cftsetrmarg{4em}
\titleformat{\section}{\Large\bfseries\raggedright}{\thesection}{0.7em}{}
\pagestyle{fancy}\fancyhf{}
\fancyhead[L]{\small\sffamily Additional UnsolvedMath statements}
\fancyhead[R]{\small Original catalogue IDs}
\fancyfoot[C]{\thepage}
\setlength{\parindent}{0pt}
\setlength{\parskip}{5pt plus 1pt}
\setlength{\emergencystretch}{3em}
\setcounter{tocdepth}{1}\setcounter{secnumdepth}{1}
\setlist[itemize]{leftmargin=3em,itemsep=4pt,topsep=4pt,parsep=0pt}
\allowdisplaybreaks
\newcommand{\N}{\mathbb{N}}
\newcommand{\Z}{\mathbb{Z}}
\newcommand{\Q}{\mathbb{Q}}
\newcommand{\R}{\mathbb{R}}
\newcommand{\C}{\mathbb{C}}
\newcommand{\F}{\mathbb{F}}
\newcommand{\A}{\mathbb{A}}
\newcommand{\PP}{\mathbb{P}}
\newcommand{\E}{\mathbb{E}}
\newcommand{\eps}{\varepsilon}
\newcommand{\dd}{\,\mathrm d}
\newcommand{\sref}[2]{\hyperlink{source:#1:#2}{[S#2]}}
\newif\ifshowcatalogue
\showcataloguetrue
\newcommand{\cataloguescope}[1]{\ifshowcatalogue
 \paragraph{Statement recorded in the supplied source.}
 #1\fi}
\begin{document}
% UnsolvedMath ID 2021; source catalogue code EP-271

\setcounter{section}{2020}

\section{Explicit Growth of Two-Seed Stanley Sequences}\label{2021}

{\small\sffamily UnsolvedMath ID 2021 · Number Theory}

\subsection{Definitions and mathematical statement}

\paragraph{Shared notation.}Unless a section says otherwise, $\N=\{1,2,\ldots\}$,
$\Z$ is the ring of integers, $\Q,\R,\C$ are the rational, real, and complex
fields, $[N]=\{1,\ldots,\lfloor N\rfloor\}$, and $\log$ is the natural
logarithm. For real functions of a parameter tending to infinity,
$f\ll g$ and $f=O(g)$ mean $|f|\leq Cg$ eventually for some fixed $C$;
$f\gg g$ reverses this comparison; $f=o(g)$ means $f/g\to0$;
$f\sim g$ means $f/g\to1$; and $f\asymp g$ means both order bounds.
A subscript on $O$ or $\ll$ specifies allowed constant dependence.
The expression $n^{o(1)}$ in an upper bound means growth smaller than
$n^\epsilon$ for each fixed $\epsilon>0$. Local definitions govern any
exception, finite-field convention, or use of the same symbol for another
object. An asymptotic claim is not automatically an effective algorithm.



The greedy choice is the least integer strictly larger than the last term that preserves absence of a nonconstant three-term progression. The inequality follows from the increasing enumeration in the source. The seed $n$ is a fixed positive integer. \sref{2021}{1} \sref{2021}{2}

\paragraph{Statement recorded in the supplied source.}
Let $A(n)=\{a_0<a_1<\cdots\}$ be the sequence defined by $a_0=0$ and $a_1=n$, and for $k\geq 1$ define $a_{k+1}$ as the least positive integer such that there is no three-term arithmetic progression in $\{a_0,\ldots,a_{k+1}\}$.
Can the $a_k$ be explicitly determined? How fast do they grow? \sref{2021}{1}

\paragraph{Residual task reported by the archive.}

The embedded review (2026-08-17) records: Determine the terms or growth for general n and prove or disprove the Odlyzko-Stanley regular/irregular dichotomy; A(4) remains a central unresolved example. \sref{2021}{1}

\subsection{Short English statement}

Starting from zero and one prescribed positive seed, how can the greedy progression-free sequence be described, and how quickly does it grow?

\subsection{Sources}

{\small

\begin{itemize}

\item[\hypertarget{source:2021:1}{S1}] ULAM.AI, \emph{UnsolvedMath}, supplied \texttt{problems.json}, record 2021 (EP-271), statement and background. The embedded literature-review date is 2026-08-17; that is the archive’s review date, not a fresh certification. Dataset landing page: \url{https://huggingface.co/datasets/ulamai/UnsolvedMath}.

\item[\hypertarget{source:2021:2}{S2}] T. F. Bloom, \emph{Erdős Problems}, Problem 271, \url{https://www.erdosproblems.com/271}. Reference imported from the supplied record; the live problem page was not independently checked for this compilation.

\item[\hypertarget{source:2021:3}{S3}] Richard A. Moy, On the Growth of the Counting Function of Stanley Sequences, Discrete Mathematics 311 (2011), 560-562, arXiv:1101.0022. \url{https://arxiv.org/abs/1101.0022}. Bibliographic information imported from the supplied record.

\item[\hypertarget{source:2021:4}{S4}] S. Lindhurst, An investigation of several interesting sets of numbers generated by the greedy algorithm. Senior thesis at Princeton University (1990). Bibliographic information imported from the supplied record.

\item[\hypertarget{source:2021:5}{S5}] A. Odlyzko and R. Stanley, Some curious sequences constructed with the greedy algorithm. Bell Laboratories internal memorandum (1978). Bibliographic information imported from the supplied record.

\end{itemize}

}

\label{end:2021}
\end{document}
